\documentclass[12pt]{article}
\usepackage[margin=2.5cm]{geometry}
\usepackage{amsmath, amssymb}
\usepackage{bm}
\usepackage{parskip}

\title{Conditional likelihood for ascertained sibships}
\author{}
\date{}

\begin{document}
\maketitle

\section*{Setup}

A sibship of size $n$ is sampled because it contains at least one affected
member.  We observe all $n$ outcomes $y_1,\ldots,y_n \in \{0,1\}$, with
$m = \sum_i y_i \geq 1$ cases.  The genotypic values $G_1,\ldots,G_n$ are
unobserved.  Under the logistic model, outcomes are conditionally independent
given genotypic values:
\[
  P(y_i = 1 \mid G_i) \;=\; \sigma(\beta_0 + G_i) \;=\;
  \frac{e^{\beta_0+G_i}}{1+e^{\beta_0+G_i}}, \qquad
  1 - \sigma(\beta_0+G_i) \;=\; \frac{1}{1+e^{\beta_0+G_i}}.
\]

\section*{Ascertainment-corrected likelihood}

The sibship is observed only if at least one member is affected.  The
probability of ascertainment given the genotypic values is
\[
  P(\text{at least one case} \mid \bm{G})
  \;=\;
  1 - \prod_{i=1}^n [1 - \sigma(\beta_0+G_i)]
  \;=\;
  1 - \prod_{i=1}^n \frac{1}{1+e^{\beta_0+G_i}}.
\]
The ascertainment-corrected likelihood is the joint outcome probability divided
by this selection probability:
\begin{equation}
  L(\beta_0, s \mid \bm{y}, \bm{G})
  \;=\;
  \frac{P(\bm{y} \mid \bm{G})}{P(\text{at least one case} \mid \bm{G})}
  \;=\;
  \frac{
    \displaystyle\prod_{i:\,y_i=1} \sigma(\beta_0+G_i)
    \;\cdot\;
    \displaystyle\prod_{i:\,y_i=0} [1-\sigma(\beta_0+G_i)]
  }{
    1 - \displaystyle\prod_{i=1}^n [1-\sigma(\beta_0+G_i)]
  }.
\end{equation}

Multiplying numerator and denominator by $\prod_i(1+e^{\beta_0+G_i})$ and
writing $e^{\beta_0+G_i}$ for the odds, the factors $(1+e^{\beta_0+G_i})^{-1}$
cancel pairwise between numerator and denominator to give:

\begin{equation}
  \boxed{
    L(\beta_0, s \mid \bm{y}, \bm{G})
    \;=\;
    \frac{
      \displaystyle\prod_{i:\,y_i=1} e^{\,\beta_0 + G_i}
    }{
      \displaystyle\prod_{i=1}^n \bigl(1 + e^{\,\beta_0 + G_i}\bigr) \;-\; 1
    }.
  }
\end{equation}

\section*{Denominator as a sum over non-empty subsets}

Expanding the product in the denominator,
\[
  \prod_{i=1}^n (1+e^{\beta_0+G_i})
  \;=\;
  1 + \sum_{\mathcal{S} \neq \emptyset}
      \prod_{i \in \mathcal{S}} e^{\,\beta_0+G_i}
  \;=\;
  1 + \sum_{k=1}^n e^{\,k\beta_0}
    \sum_{\mathcal{S}:\,|\mathcal{S}|=k}
    e^{\,\sum_{i\in\mathcal{S}} G_i},
\]
so the denominator is
\[
  \sum_{k=1}^n e^{\,k\beta_0}
  \sum_{\mathcal{S}:\,|\mathcal{S}|=k} e^{\,\sum_{i\in\mathcal{S}} G_i}.
\]
This is a weighted sum over all $2^n - 1$ non-empty subsets $\mathcal{S}$,
each contributing $e^{|\mathcal{S}|\beta_0 + \sum_{i\in\mathcal{S}} G_i}$.
Because different subset sizes $k$ enter with weights $e^{k\beta_0}$ that
depend on $\beta_0$, \textbf{the intercept does not cancel}: the likelihood
retains information about the prevalence $K = \sigma(\beta_0)$.

\section*{Special cases}

\subsection*{Discordant pair ($n=2$, $m=1$)}

Without loss of generality let $y_1=1$, $y_2=0$:
\[
  L \;=\;
  \frac{e^{\beta_0+G_1}}{(1+e^{\beta_0+G_1})(1+e^{\beta_0+G_2}) - 1}
  \;=\;
  \frac{e^{\beta_0+G_1}}{e^{\beta_0+G_1} + e^{\beta_0+G_2} + e^{2\beta_0+G_1+G_2}}.
\]
For rare disease ($K \ll 1$, $\beta_0 \to -\infty$) the term
$e^{2\beta_0+G_1+G_2}$ is negligible, and
\[
  L \;\approx\;
  \frac{e^{G_1}}{e^{G_1}+e^{G_2}} \;=\; \sigma(G_1 - G_2),
\]
which is identical to the conditional logistic likelihood for a discordant pair.

\subsection*{Concordant-affected pair ($n=2$, $m=2$)}

\[
  L \;=\;
  \frac{e^{2\beta_0+G_1+G_2}}{e^{\beta_0+G_1} + e^{\beta_0+G_2} + e^{2\beta_0+G_1+G_2}}.
\]
For rare disease this is approximately $e^{\beta_0}(e^{G_1}+e^{G_2})^{-1}$, so
the likelihood shrinks as $K \to 0$: a concordant-affected pair is increasingly
surprising for a rare disease, and this information is retained.  The
conditional logistic likelihood for a concordant pair is 1 (trivially), so
\textbf{concordant-affected sibships carry information here but not in the
conditional logistic likelihood}.

\section*{Comparison with the conditional logistic likelihood}

The conditional logistic likelihood conditions on the observed stratum total
$m = \sum_i y_i$:
\[
  P(\bm{y} \mid m,\, \bm{G})
  \;=\;
  \frac{e^{\,\sum_{i:\,y_i=1} G_i}}
       {\displaystyle\sum_{\mathcal{S}:\,|\mathcal{S}|=m}
        e^{\,\sum_{i\in\mathcal{S}} G_i}}.
\]
Compared with the ascertainment-corrected likelihood:
\begin{itemize}
  \item The denominator sums over \emph{all $2^n-1$ non-empty subsets}
        rather than only the $\binom{n}{m}$ subsets of the observed size $m$.
  \item $\beta_0$ does not cancel, because the weights $e^{k\beta_0}$ differ
        across subset sizes.
  \item Concordant-affected sibships ($m=n$) contribute to the likelihood,
        whereas their contribution to the conditional logistic likelihood is
        trivially~1 (a single term in numerator and denominator).
  \item For rare disease and discordant strata, the two likelihoods agree to
        first order in $K$.
\end{itemize}
The conditional logistic likelihood can be recovered by factoring the
ascertainment-corrected likelihood as
\[
  L(\beta_0, s \mid \bm{y}, \bm{G})
  \;=\;
  P(\bm{y} \mid m, \bm{G})
  \;\times\;
  \frac{\displaystyle\sum_{\mathcal{S}:\,|\mathcal{S}|=m}
        e^{\,\beta_0 m + \sum_{i\in\mathcal{S}} G_i}}
       {\displaystyle\sum_{k=1}^n e^{k\beta_0}
        \sum_{\mathcal{S}:\,|\mathcal{S}|=k}
        e^{\,\sum_{i\in\mathcal{S}} G_i}},
\]
where the second factor is the probability that the observed $m$ is the
stratum total, given that the stratum has at least one case.  The conditional
logistic likelihood discards this second factor and its information about
$\beta_0$.

\end{document}
